\section{Post-Baseline Performance and Constant-Time Work}
\label{sec:post-baseline}

With baseline interoperability achieved (\S\ref{sec:status}), the
implementation moves into a post-baseline phase: performance
optimization and constant-time hardening built on top of the
byte-equality foundation.  Each change is interop-safe (KAT bytes
preserved) unless explicitly noted under
\S\ref{sec:tried-rejected}.  The field-arithmetic and elliptic-curve
layers are already constant-time (the Montgomery ladder runs a fixed
bit count; \texttt{subtle::ConstantTimeEq} guards secret comparisons);
the quaternion-arithmetic layer is the primary remaining variable-time
surface, and the entries below describe its progressive hardening.

The remainder of this section is two running logs --- one of
experiments that landed, one of experiments that did not.  Recording
the second is as important as recording the first: each entry that
names a tried-and-rejected approach spares the next reader the same
dead end.

\subsection{Landed work}
\label{sec:perf-ct}

Measurements below are the deterministic instruction-count delta
(Valgrind via \texttt{gungraun}) relative to \texttt{main} at the time
the work landed, on the operation indicated.  Wall-clock benchmarks
(\texttt{divan}) are tracked separately because they are sensitive to
runner conditions.

\paragraph{MLLL replacing dual-sum-dual at the response-phase
intersection.}
The response phase computes the intersection
$I_{\mathrm{sk}} \cap \bar{I}_{\mathrm{com}}$ via the dual-sum-dual
identity, which required a working width near $500$~limbs to absorb
HNF coefficient blow-up.  We replaced the HNF reduction with the
modified-LLL (MLLL) of Kim, Lee, and Yoo~\cite{KLY26}, dropping the
working width to $128$~limbs at the main intersection site and
$64$~limbs at the smaller auxiliary one.  Cost scales $O(W^2)$, so
the narrowing is large.  Deterministic measurement: $-11.8\%$ on
signing, output byte-identical.  A follow-up narrowed the widths
further to the $\leq 110$-limb bound proven by Kim et al.\ (ePrint
2025/1649), giving $110$~limbs at the main site and $48$ at the
auxiliary site.

\paragraph{Constant-time modular exponentiation in primality testing.}
\texttt{MontReducer::pow} (Miller--Rabin during the norm-equation
prime search of \texttt{represent\_integer}, where the exponent
$d = (n-1)/2^s$ is secret-derived) was made constant-time: the window
table is read with a scatter-gather \texttt{ct\_select} so the access
pattern does not leak the exponent nibbles, and the final Montgomery
subtracts in \texttt{mul}, \texttt{reduce\_wide}, and
\texttt{compute\_r\_squared} are now branchless.  Output is
identical.  Deterministic measurement: $-1.01\%$ on key
generation, $-0.10\%$ on signing.  Verification untouched.

\paragraph{Hardware AES on aarch64.}
The AES-CTR-DRBG that drives every random draw in key generation and
signing (roughly $3000$ block encryptions per signature) was running
the \texttt{aes} crate's software fixslice backend on aarch64.  The
crate gates its ARMv8 Cryptography Extensions backend behind a
\texttt{cfg(aes\_armv8)} that the build was not setting, whereas the
x86\_64 AES-NI backend is compiled unconditionally and selected at
runtime.  Setting the cfg on the aarch64 target tables lets
\texttt{cpufeatures} runtime-select the hardware backend, with the
software path preserved as a fallback on CPUs lacking
\texttt{FEAT\_AES}.  Hardware AES is constant-time and bit-identical to
the fixslice output, so all KATs stay byte-identical.  Wall-clock A/B
(Apple M4): $-12.0\%$ on key generation, $-3.7\%$ on signing.  Key
generation gains more because the norm-equation prime search issues far
more DRBG draws per operation than signing does.  A time-profile of the
hardware-AES build shows the residual AES cost (about $1.25\%$ of key
generation, less of signing) is dominated by the AES key schedule, not
block encryption.  The DRBG was also switched from the full
\texttt{aes::Aes256} cipher to the encryption-only \texttt{Aes256Enc},
since CTR mode never decrypts; this drops the wasted decryption key
schedule and is the correct cipher type for the construction.

\paragraph{One key schedule per DRBG generate cycle.}
The key-schedule cost above is one AES-256 expansion per
\texttt{CTR\_DRBG\_Generate}: the \texttt{Key} changes on every
\texttt{Update}, so one expansion per cycle is mandated by SP 800-90A
and is not byte-identically removable.  The implementation was doing
\emph{two}, though: \texttt{randombytes} expanded the key for its
generate loop, then the trailing \texttt{Update} (which encrypts under
the same, not-yet-changed \texttt{Key}) expanded it again.  Sharing the
one schedule across both loops halves the per-cycle expansion; the
generate loop was also switched from one-block-at-a-time
\texttt{encrypt\_block} to batched \texttt{encrypt\_blocks} so the
backend's parallel pipeline is fed.  Output is byte-identical (all KATs
green).  Profiled on an Apple M4, the key-schedule self-time in key
generation fell from about $1.05\%$ to $0.50\%$ of the total, with the
generate-loop encryption falling below the sampling threshold.

\paragraph{Leading-zero skipping in the wide integer multiply.}
The quaternion intersection and Hermite-normal-form arithmetic runs
multi-precision multiplication at integer widths sized for the worst
case (the proven $110$-limb bound of Kim et al.\ at the main
intersection, $48$ and $30$ limbs elsewhere), but a profile of the
actual operands shows the typical value occupies under half those
limbs, so the truncated schoolbook multiply was spending roughly half
its time multiplying zero limbs.  The multiply now bounds its loops by
the operands' effective limb lengths and propagates the final carry,
producing byte-identical limbs.  This matches the C reference's bignum
library, whose multiply already operates only on the significant limbs;
the fixed-width grid was an artifact of our worst-case storage width.
Wall-clock A/B (Apple M4): $-44.7\%$ on signing, $-14.4\%$ on key
generation, verification unchanged; all KATs and the multi-precision
differential oracle stay byte-identical.  The win is large because the
wide multiply dominates signing once MLLL replaced the dual-sum-dual
intersection.  Constant-time scope: the skip is restricted to the wide
widths, so it is reached only by the variable-time quaternion-lattice
arithmetic; the narrow-width multiplies that every constant-time path
relies on (the secret-torsion endomorphism action; the field layer uses
none; primality's modular exponentiation has its own routines) keep the
fixed-width path and are unaffected.  Like the reference's bignum
library this path is variable-time in operand magnitude; making it
constant-time folds into the quaternion layer's eventual constant-time
pass, which reworks these multiplies.

\subsection{Tried approaches that did not land}
\label{sec:tried-rejected}

Each entry below names what was tried and the reason it did not
land.

\paragraph{Cost-optimal isogeny strategy for the verify-side 4-isogeny
chain.}
The verify-side $2^e$-isogeny chain in
\texttt{src/curves/isogeny/mod.rs} --- decomposed into
$\lfloor e/2 \rfloor$ 4-isogenies --- was changed from runtime
balanced halving to a compile-time-precomputed cost-optimal De
Feo--Jao--Plut strategy.  Output is byte-identical.
Deterministic measurement:
$-0.04\%$ on verification, $0$ on key generation and
signing --- no measurable gain because verification's hot path
is $\mathbb{F}_{p^2}$ and elliptic-curve arithmetic, not the strategy
walker, at our chain length ($e \leq 126$).

\paragraph{Cost-optimal strategy on the $(2,2)$-theta chain (surfaces).}
The same DP strategy was applied to Phase~1 of the $(2,2)$-theta chain
in \texttt{src/surfaces/mod.rs}.  $166$ of $707$ lib tests failed,
including the majority of signing KATs.  Root cause: Phase~1 uses the
non-standard Jacobian doubling \texttt{double\_for\_theta} (with
$Z' = 2yz$) deliberately, so the gluing step receives the specific
projective representative the C reference produces.  Reordering the
doublings within the strategy walker changes the cumulative $Z$
scaling, which the gluing matrix multiplies by, which then changes the
produced signature bytes --- even though the composed isogeny is
mathematically unchanged.  Cleared to the non-interop pass.

\paragraph{Batch inversion in \texttt{cross\_pairings} final
exponentiation.}
The five independent $\mathrm{den}^{-1}$ inversions in the
\texttt{cross\_pairings} final-exponentiation loop were folded into a
single inversion via Montgomery's trick.  Output is identical and
every KAT byte-identical.  However deterministic measurement showed
a $+0.05\%$ regression overall and $+0.62\%$ on verification: the
two-pass refactor's overhead
(extra $\mathbb{F}_{p^2}$ array storage, a likely inlining-decision
change in a now-larger function) exceeded the savings from fewer
inversions, given how rarely \texttt{cross\_pairings} is invoked per
signing operation.

\paragraph{Substituting MLLL into ideal multiplication.}
Replacing the HNF step in \texttt{Lattice::product} with MLLL on the
same theory as the intersection win.  In practice the multiplication
runs at small widths where MLLL's fixed overhead --- a $16 \times 16$
Gram matrix and the iterative reduction --- exceeds the modest HNF
blow-up it removes.  Result: roughly $1.9\times$ regression on
key generation in a controlled A/B.

\paragraph{Input-LLL preprocessing of the intersection.}
The input-LLL preprocessing of~\cite{KLY26} reduces each input basis
before dualization, intending to bound the duals' cubic coefficient
growth.  Our existing \texttt{l2\_reduce} routine minimizes the
reduced-norm form $a^2 + b^2 + p\,(c^2 + d^2)$, which is the natural
norm in $B_{p,\infty}$ but is dominated by the $p\,(c^2 + d^2)$ term;
the coordinate magnitudes that drive the Euclidean blow-up therefore
remain unbounded under our reduction.  The intersection width did not
fall below $128$~limbs and the experiment was shelved pending a
Euclidean reduction step (the constant-time dim-4 LLL of Hanyecz et
al., ePrint 2025/027, is the intended substitute).

\paragraph{Hanyecz constant-time dim-4 BKZ-2 as preprocessing in front
of the MLLL intersection.}
Hanyecz et al.\ (ePrint 2025/027) give a constant-time dim-4 BKZ
with block size $2$ as the Euclidean-norm replacement for
\texttt{l2\_reduce} on secret-derived inputs.  Their five algorithms
port to \texttt{src/quaternions/lattice/bkz2.rs} on branch
\texttt{hanyecz-euclidean-input-lll}; invoked in front of
\texttt{compact\_intersection} at the two response-phase call sites
of \S\ref{sec:perf-ct}, output stays byte-identical with the C
reference (canonical-HNF depends only on the inputs'
$\mathbb{Z}$-modules, not their basis representatives).

The integration did not land.  MLLL already runs at the proven
width bound of~\cite{KLY26}, which is a property of the lattice and
not of any basis representation, so BKZ-2 preprocessing cannot
narrow it further; in our measurements the preprocessing added cost
at every parameter setting tried and regressed signing wall-clock
by roughly an order of magnitude.  Separately, the paper's own
proposed use of BKZ-2 --- substituting for \texttt{l2\_reduce} on
secret-derived inputs --- is KAT-breaking, since the surviving
\texttt{l2\_reduce} call sites are followed by basis-dependent
downstream logic.  Capturing the paper-claimed CT win therefore
requires a KAT re-baseline; it bundles with the Cornacchia fix
(ePrint 2025/2192) and the fast-primality items (\S\ref{sec:next})
into a single non-interop release.  The four BKZ-2 algorithm
commits remain on \texttt{hanyecz-euclidean-input-lll} as a
self-contained library; the integration commit is shelved.

\paragraph{Whole-program LTO with a single codegen unit.}
A \texttt{[profile.release]} with \texttt{lto = "fat"} and
\texttt{codegen-units = 1} was tried on the theory that cross-crate
inlining of the \texttt{aes}, \texttt{subtle}, and field kernels would
help.  It regressed both hot operations relative to the default
sixteen-codegen-unit build: wall-clock A/B (Apple M4) of $+6\%$ on
signing and $+5\%$ on key generation.  The single codegen unit appears
to make worse register-allocation and inlining choices in the MLLL and
lattice-product inner loops than the per-unit default.  Output is
byte-identical; the regression is pure codegen.  Not adopted.

\paragraph{Byte-identical micro-optimizations with no wall-clock
effect.}
Three byte-identical changes measured neutral (within run-to-run noise)
on an Apple M4 and were not adopted as performance work: (a) a
stack-allocated $16$-element array in place of the \texttt{Vec}
accumulator in \texttt{Lattice::product} and its modulus variant,
removing the heap reallocations of the $16$-column product but not
moving wall-clock; (b) fusing the doubling and conditional-subtract
limb passes in the $R^2$ Montgomery-setup loop; and
(c) strength-reducing the multiply-by-$p$ in the MLLL Gram computation
to $(x \ll 250) + (x \ll 248) - x$ (using $p = 5\cdot 2^{248} - 1$),
which at the $110$-limb response-phase width is no faster than the
length-aware schoolbook multiply by the four-limb constant it
replaced.

\subsection{Next steps}
\label{sec:next}

\begin{enumerate}
  \item Costello--Hisil 4-isogeny formula and Renes square-root-free
    2-isogeny formula in the elliptic factors (verify-side chain),
    composed underneath the cost-optimal strategy already in place.
  \item Leroux CRT intersection~\cite{basso2025}, replacing the
    dual-sum-dual identity entirely; the roadmap claims
    signatures-identical, but the surfaces $(2,2)$-theta result above
    means this should be verified with a small prototype before the
    full substitution.
  \item Continue the constant-time hardening sequence
    (\S\ref{sec:ct}).  The remaining variable-time surfaces are
    \texttt{MontReducer::pow}'s base reduction
    (\texttt{mag\_div\_rem}), Cornacchia in the norm-equation search
    (ePrint 2025/2192, the highest-severity item --- poly-time key
    recovery from the input leak), and \texttt{l2\_reduce} on
    secret-derived input (where the Hanyecz BKZ-2 substitution above
    is the intended fix, bundled with the other KAT-breaking CT items
    into a single non-interop release).
  \item In-place / out-parameter refactor of the hot wide-\texttt{BigInt}
    lattice operations to remove struct-copy overhead
    (\texttt{memcpy} is $\sim\!13.5\%$ of signing).  The arithmetic
    returns values by move --- \texttt{Lattice::product},
    \texttt{Vector}/\texttt{Matrix} operations, and the \texttt{BigInt}
    operators all return new wide structs (a $4\times 4$
    \texttt{Matrix} of $30$-limb integers is $\sim\!3.8$~KB), copied on
    every compose in the HNF / intersection loops.  The C reference
    avoids this with an out-parameter, const-reference convention
    (\texttt{op(out, a, b)}, e.g.\ \texttt{ibz\_vec\_4\_add}); adopting
    in-place variants for the hot wide operations is the corresponding
    structural lever, and is interop- and constant-time-safe.  Unlike
    the assembly attempts (\S\ref{sec:backends}), it is a portable,
    all-architecture change.
\end{enumerate}
